% ================= 3. SYSTEM ANALYSIS & REQUIREMENTS =================
\section{System Analysis \& Requirements}

\subsection{Scope of Work \& Use Cases}
\label{sec:scope}

% Define the boundaries of this design effort and the concrete scenarios the
% system must support. Reference the four target behaviors from Sec.~\ref{sec:objectives}
% (edge balance, corner balance, controlled rotation, walking) as the primary use cases.

\subsection{Actuation Trade-off --- Why Reaction Wheels}
\label{sec:tradeoff}

% Justify momentum-exchange actuation over alternatives (e.g. control moment
% gyros, external actuation). Build on the ADCS framing of Sec.~\ref{sec:application}
% and the RW-vs-CMG distinction \cite{votel2012cmgvsrw, tensortech2026rwcmg}.

\subsection{Requirements \& Objectives Traceability}
\label{sec:requirements}

Each project objective (Section~\ref{sec:objectives}) imposes a required capability,
which in turn is confirmed by a specific test in the validation plan
(Section~\ref{sec:testing}). The mapping below keeps every requirement traceable
from motivation to verification.

\begin{table}[h]
\centering
\caption{Objective-to-requirement-to-verification traceability.}
\label{tab:traceability}
\begin{tabular}{lll}
\toprule
Objective & Required capability & Verified by \\
\midrule
Edge balance        & Regulation about a line-contact equilibrium & Sec.~\ref{sec:testing} \\
Corner balance      & Regulation about a point-contact equilibrium, all axes & Sec.~\ref{sec:testing} \\
Controlled rotation & Reference tracking / slew between equilibria & Sec.~\ref{sec:testing} \\
Walking (stretch)   & Chained jump-up + re-stabilization & Sec.~\ref{sec:testing} \\
\bottomrule
\end{tabular}
\end{table}

\subsection{Coordinate Frames \& Conventions}
\label{sec:frames}

% Define the body frame, world/inertial frame, and the sign conventions for
% tilt angle and wheel angle used consistently throughout the modeling section.
